\section{Weight flow, information geometry and Legendre duality}
\label{sec:convex-refinement}

Interval deformation admits an exact variational
description. Let \(d_j=K_j/Q>0\), \(\sum_jd_j=1\), and let \(u_j\) be the
coordinates of the direction \eqref{eq:u-main}. For \(s\in\RR\),
define
\begin{equation}
 Z(s)=\sum_{j=1}^{12}d_j e^{su_j},\qquad
 \psi(s)=\log Z(s),\qquad
 d_j(s)=\frac{d_j e^{su_j}}{Z(s)}.
 \label{eq:convex-partition}
\end{equation}
All exponential arguments are dimensionless. The parameter
\(s\) is the coordinate of this geometric deformation. Assigning
a physical duration belongs to a declared temporal reader.

\begin{theorem}[Strict convexity and dual coordinate]
The function \(\psi\) is analytic and strictly convex on \(\RR\).
If \(u_- =\min_j u_j\) and \(u_+=\max_j u_j\), the map
\(v=\psi'(s)\) is a diffeomorphism from \(\RR\) onto \((u_-,u_+)\).
For each \(v\) in that interval there is a unique coordinate \(s(v)\), and
\[
 \psi^*(v)=s(v)v-\psi(s(v)),\qquad
 (\psi^*)'(v)=s(v),\qquad
 (\psi^*)''(v)=\frac1{\psi''(s(v))}>0.
\]
\end{theorem}
\begin{proof}
The finite sums allow termwise differentiation. This gives
\begin{align}
 \psi'(s)&=\sum_jd_j(s)u_j=:\overline u(s),\notag\\
 \psi''(s)&=\sum_jd_j(s)(u_j-\overline u(s))^2.
 \label{eq:variance}
\end{align}
The weights are strictly positive and the \(u_j\) are not all equal;
the variance is therefore strictly positive. As \(s\) tends to
\(+\infty\), dividing numerator and denominator of \(\psi'\) by
\(e^{su_+}\) shows that only indices with \(u_j=u_+\) survive.
Thus \(\psi'(s)\to u_+\). The argument with \(u_-\) gives the opposite
limit. Monotonicity and the inverse function theorem provide
smooth bijectivity. The function \(sv-\psi(s)\) has its unique maximum
at \(s(v)\); differentiating the maximum value proves the two dual
identities.
\end{proof}

The evolution of each interval is
\begin{equation}
 \dot d_j(s)=d_j(s)\bigl(u_j-\overline u(s)\bigr).
 \label{eq:replicator}
\end{equation}
The derivatives sum to zero. Total length therefore remains
one while the internal lengths are redistributed. This conservation
of global scale allows the change in cross-ratios to be attributed to a
deformation of positions rather than a common dilation.

On the open simplex of weights, the metric
\(g_d(\xi,\zeta)=\sum_j\xi_j\zeta_j/d_j\), on tangent vectors
of sum zero, is positive. The linear function
\(F(d)=\sum_jd_ju_j\) has, with respect to this metric, exactly the
field in \eqref{eq:replicator} as its gradient, since
\[
 g_d\bigl(d_j(u_j-\overline u),\xi\bigr)
 =\sum_j(u_j-\overline u)\xi_j=\sum_ju_j\xi_j=dF(\xi).
\]
Consequently \(dF/ds=\psi''(s)>0\) along the flow. The
ascending character corresponds to this specific function and metric.

\begin{theorem}[Variational invariance under inherited subdivision]
For each interval \(j\), let \(r_{j,c}>0\),
\(\sum_c r_{j,c}=1\), be the relative weights of a fixed subdivision.
Assign to the subintervals
\[
 d_{j,c}=r_{j,c}d_j,\qquad u_{j,c}=u_j.
\]
The partition function, its logarithm, all its derivatives and the Legendre
transform agree exactly with those of the predecessor partition.
Furthermore, the flow commutes with summation over subintervals:
\[
 \sum_c d_{j,c}(s)=d_j(s),\qquad
 d_{j,c}(s)=r_{j,c}d_j(s).
\]
\end{theorem}
\begin{proof}
Grouping the finite sum gives
\[
 Z_{\rm ref}(s)=\sum_{j,c}r_{j,c}d_j e^{su_j}
 =\sum_j d_j e^{su_j}=Z(s).
\]
All conclusions follow from this identity and from dividing each
weight by the same \(Z(s)\). The same argument holds after any
finite number of subdivisions. For a countable subdivision, the sum of
positive terms converges to the same value when grouped by predecessor intervals.
\end{proof}

The inherited-direction hypothesis has verifiable content. If an
interval with direction \(a\) is divided into two halves with directions
\(a+b\) and \(a-b\), its contribution changes from
\(d e^{sa}\) to \(d e^{sa}\cosh(sb)\). For \(b\ne0\), equality
of partition functions fails even though the total weight at \(s=0\) is
the same. Geometric and variational conservation thus requires preserving
the relevant directional information as well. The proved identity
applies to inherited subdivision; its transport to a TPK update
that changes those directions is determined by composing the reader of that
update.

There is a regular Hamiltonian formulation of the dual coordinate. In the
cotangent plane with coordinates \((q,p)\), take
\(H(q,p)=\psi(p)\). The equations are
\(\dot q=\psi'(p)\), \(\dot p=0\). For velocities
\(\dot q\in(u_-,u_+)\), the Legendre transform gives
\[
 L(q,\dot q)=\psi^*(\dot q),\qquad
 \frac{\partial^2L}{\partial\dot q^2}>0.
\]
This system variationally realizes the potential of the interval flow.
Its equation \(\dot p=0\) shows that Hamiltonian time differs
from the parameter \(s\) traversing the collection \eqref{eq:convex-partition}.
Evaluation of \(H\) and \(L\) is preserved under the subdivision of the
theorem, because \(\psi\) and \(\psi^*\) remain identical.

In turn, any flow \(\dot d=F(d)\) admits, in a local chart
of the simplex, a cotangent lift with
\(\mathscr H(d,p)=p\cdot F(d)\). This second construction reproduces
the flow \eqref{eq:replicator} itself on its base, but its Hessian with respect
to \(p\) is zero. The two realizations answer different questions:
one produces a regular convex duality and the other transports dynamics
on a given base. This distinction preserves the meaning of Hamiltonian and
Lagrangian when comparing this geometry with the physical readers of the corpus.
